\honest{My Childhood Death Spiral}{Eliezer Yudkowsky}{2008}

My parents always played down intelligence and played up experience: ``maturity and wisdom is more important than intelligence.'' When I argued against Judaism, I was told that ``Logic has limits'' and that when I was older experience would show me the truth of Judaism. I did not try again. If this was meant to make me value intelligence above all, it was the most successful reverse psychology I know of. But my parents are not that cunning, and the results were not good.

For a long time I thought the moral was that raw intelligence beats experience. Later I saw that an eleven-year-old could not have out-thought adult parents in a fair contest. My SAT scores were high for my age but would not have beaten theirs. They lost because of ``dysrationalia'': they used their intelligence to defeat itself.

The moral I drew as a child was that anyone who plays down intelligence does not understand it. My intelligence had shaped every part of my life and personality. Does self-awareness have nothing to do with wisdom or goodness? Modeling yourself takes intelligence, if only enough to learn evolutionary psychology. Intelligence is ``the unfairest'' of the cards we are dealt, more unfair than wealth, health, home country or one's happiness set-point. Saying ``Intelligence isn't as important as X'' is a way of turning from that unfairness. It tempts those dealt poor cards and those dealt good ones, as playing down money tempts both the poor and the rich. \nb{The postscript says that the author now thinks first of death and old age as the poor hands dealt to humans.}

But I was a transhumanist, so low intelligence was a problem to be fixed, even if it took my whole life. ``The strong exist to serve the weak,'' I wrote, ``and can only discharge that duty by making others equally strong.'' I was reacting against the Randian and Nietzschean trends in science fiction and went too far the other way. No one exists only to serve. But I tried, and I do not regret it. Everyone needed to be smarter. I was careful to include myself, to claim no moral superiority, and to draw no line between myself and others. \nb{The author's archived essay ``Staring into the Singularity'' (dated 1996 and 1999) does say ``I know, in a dim way, just how dumb I am.''}

Having read my science fiction, I avoided the obvious traps, and did not fail in any obvious way. In 1996 I met the idea of the Singularity. It was no thunderbolt. It simply seemed obvious that smarter-than-human intelligence would change the future more deeply than any material science, including the nanotechnology I had aimed at when I was eleven, which would only be a tool of intelligence. I knew at once that I would spend my life creating the Singularity. Intelligence was even more powerful than I had thought. This was a ``happy death spiral'': it led me to false happy beliefs about intelligence. Perhaps the line is where I came to believe that the speed of light would surely be no barrier to a superintelligence (not unthinkable, but I would not bet on it). \nb{Yudkowsky's early texts that the editor could read do not state that belief, though ``Staring into the Singularity'' mocks people who set ``hard limits on Powers.''}

The real wrong turn came when someone said that intelligence has nothing to do with being good. My own ethics were built with intelligence; just try explaining the Prisoner's Dilemma to a chimpanzee. So I concluded that superintelligence ``would necessarily imply supermorality.'' \nb{The early writings bear this out. ``Staring into the Singularity'': ``I am reasonably (95\%) certain that the Powers will be ethical.'' ``The Singularitarian Principles'' (2000, revised 2001): ``There was a point where I was sure that superintelligent meant super-ethical (probably true).'' The case is the kind of real example that ``Resist the Happy Death Spiral'' lacked.}

``Parents do all the things they tell their children not to do, which is how they know not to do them.'' To be continued.

In a postscript: what matters about intelligence, from the point of view of fun, is that it should increase over time instead of declining as it does now. I no longer work so hard at downplaying my intelligence, because that only calls attention to it. Intelligence is still ``the lever that lifts worlds,'' but it is less mysterious to me now; I see it as something inside physics. Superintelligences may travel faster than light if the true laws of physics permit it, ``and if not, then not.''
